% Transcription — fonds Alexandre Grothendieck, shelfmark 116, batch 1 (pages 1-10).
% Established from the facsimile archives/batches/116/batch-01.pdf (a mirror of
% https://grothendieck.umontpellier.fr/116.pdf), per the skill
% transcribe-grothendieck. The folder is ten pages and this is its only batch.
%
% Seven of the ten leaves carry his hand and all seven are transcribed. Page 1
% is the kraft wrapper the notes were folded into; it carries two words of his,
% « Mapping cone », alone in the top right corner, and nothing else — so it
% takes no \page{} and its words become the first section title, which is also
% where the inventory's unbracketed title comes from. Page 2 is the other face
% of that same wrapper and does carry mathematics, so it is transcribed. Pages
% 3 to 7 are the run proper, in ink on white leaves, under a title of his own:
% « Relations entre morphismes de topos et recollements de topos ». Page 9 is a
% separate leaf on a different question and is given its own section. Page 8 is
% blank (only the show-through of page 9 is visible) and page 10 is an
% administrative notice of the Université des Sciences et Techniques du
% Languedoc dated 19 June 1981 — presumably how the archivists arrived at
% « [à partir de 1981] », though nothing says so. Both are absent; the gaps in
% the page numbers are the record.
%
% The subject is one construction taken from three sides. A morphism of topoi
% f : X_0 --> X_1 is read as a topos fibred over Delta_1, and its total topos
% X = C(f) — he calls it the mapping-cone — has for sheaves the systems
% (F_0, F_1, u) with u : f^*(F_1) --> F_0. In X, X_0 is the open part and X_1
% the closed complement, the closed inclusion is strictly aspherical, and the
% relative cohomology sequence is rewritten as an H^*(X_1 mod X_0). Pages 5 to
% 7 run the construction backwards: from a topos with an open sub-topos, or
% equivalently from a morphism to Top(Delta_1); then from a triple
% (X_0, X_1, phi) with phi : F(X_0) --> F(X_1), asking what makes the two
% inclusions weak equivalences; and finally, in the case X = Top(X) with X a
% category, reducing the datum of phi to a bifunctor
% X_0^op x X_1 --> (Ens).
%
% Three things about the writing, stated once here rather than page by page.
% He replaces « asphérique » by « équivalence faible » by interlinear
% substitution, and does it twice — page 3 and page 5 — the first word being
% struck each time; both are transcribed as he left them. He writes the sign
% « = » for the word « même » (pages 4 and 6). And the left edge of page 3
% carries two blocks written in biais at about 40-45°, in a faster hand than
% the body; they are given as far as they go, which is not far, and the rest
% is marked rather than filled.
%
% One variance of substance is flagged where it occurs and not repaired: the
% adjoint f^* of phi is called « à gauche » on pages 4 and 7 and « à droite »
% on page 6.
%
% Pass: Opus 5 (claude-opus-5), 2026-09-12 — first pass, unchecked against the pages by a human.
\documentclass[11pt,a4paper]{article}
% Le préambule commun à toutes les transcriptions du fonds Grothendieck.
%
% Il tient en une idée : **l'incertitude se compose, elle ne se résout pas.**
% Une transcription qui lisse un mot douteux a détruit la seule chose pour
% laquelle on la faisait — un lecteur doit pouvoir savoir, à chaque endroit,
% ce qui était sur le feuillet et ce qui a été ajouté. Les six macros qui
% suivent sont donc l'appareil critique, et non de la décoration.
%
% Les mêmes commandes sont reconnues par `scripts/render.mjs`, qui produit la
% vue de lecture du site. Une macro ajoutée ici doit l'être aussi là-bas,
% faute de quoi le rendu échouera bruyamment — ce qui est voulu : un rendu
% silencieusement faux est pire qu'un rendu absent.

\ProvidesPackage{grothendieck}[2026/08/08 Transcriptions du fonds Grothendieck]

\usepackage{amsmath,amssymb,amsthm}
\usepackage{mathrsfs}
\usepackage{stmaryrd}
% Les diagrammes commutatifs sont partout dans ce fonds : c'est la langue
% même de la géométrie algébrique de Grothendieck, et une transcription qui
% les rendrait en prose ne transcrirait rien.
\usepackage{tikz-cd}
% Français par défaut — c'est la langue des feuillets — mais l'anglais est
% chargé aussi : la traduction et le résumé sont des documents anglais, et la
% typographie française (espaces avant les deux-points) y serait une faute.
% Ces documents basculent par \selectlanguage{english} en tête de corps.
\usepackage[english,french]{babel}
\usepackage{fontspec}
\usepackage{geometry}
\geometry{a4paper,margin=25mm,marginparwidth=28mm}
\usepackage{marginnote}
\usepackage{xcolor}
% ulem, not soul. `soul`'s \st and \ul are built on a character-by-character
% scan that XeLaTeX's font machinery breaks: the document fails with an
% undefined control sequence rather than degrading. ulem does the same two jobs
% and is Unicode-engine safe. `normalem` stops it hijacking \emph, which the
% transcriptions use for Grothendieck's own emphasis.
\usepackage[normalem]{ulem}
\usepackage{hyperref}
\hypersetup{colorlinks=true,linkcolor=black,urlcolor=black,pdfborder={0 0 0}}

% --- Métadonnées du lot -----------------------------------------------------
% Reprises telles quelles par le script de rendu, qui les affiche en tête de
% la vue de lecture. Elles fixent ce que le fichier prétend couvrir : sans
% elles, une transcription flotte sans point d'attache dans un fonds de seize
% mille feuillets.

\newcommand{\folder}[1]{\gdef\grothFolder{#1}}
\newcommand{\batch}[1]{\gdef\grothBatch{#1}}
\newcommand{\leaves}[2]{\gdef\grothFirst{#1}\gdef\grothLast{#2}}
\newcommand{\foldertitle}[1]{\gdef\grothFoldertitle{#1}}
\gdef\grothFolder{?}\gdef\grothBatch{?}\gdef\grothFirst{?}\gdef\grothLast{?}\gdef\grothFoldertitle{}

% --- Le résumé --------------------------------------------------------------
%
% Il ouvre la lecture modernisée, sous le titre « Résumé », et il est composé
% en petit. Ce n'est pas un ornement : c'est le seul passage du document qui
% ne soit pas des mathématiques, et un lecteur doit pouvoir le reconnaître
% avant de l'avoir lu — puis le sauter s'il connaît déjà le sujet, ou s'y
% arrêter s'il ne le connaît pas. Le filet qui le suit marque la reprise du
% corps.

\newenvironment{resume}
  {\small\setlength{\parskip}{0.45em}}
  {\par\medskip\hrule\medskip}

% --- Mots-clés ---------------------------------------------------------------
%
% En anglais, en clôture du résumé de la lecture modernisée : le vocabulaire
% moderne sous lequel chercher ce que les feuillets construisent. Le manifeste
% du site (`npm run manifest`) les extrait pour étiqueter la cote entière —
% cette ligne est la source unique des tags, il n'y a pas de fichier à côté.
\newcommand{\keywords}[1]{\par\smallskip\noindent
  {\footnotesize\textsc{Keywords} --- \textit{#1}}\par}

% --- L'appareil critique ----------------------------------------------------

% Le numéro de feuillet, en marge. C'est l'ancre qui permet de régler un
% désaccord sur une formule en regardant *un* feuillet plutôt qu'une chemise
% de six cents. Le site s'en sert aussi pour tourner les pages du fac-similé
% quand on fait défiler la transcription.
\newcommand{\leaf}[1]{%
  \marginnote{\footnotesize\color{gray!70}\textsf{#1}}%
  \label{leaf:#1}\ignorespaces}

% Illisible : on ne devine pas. Un mot inventé qui se lit comme les autres est
% le pire résultat possible.
\newcommand{\ill}{\textcolor{red!60!black}{[\,\ldots\,]}}

% Lecture douteuse : le mot est donné, et signalé comme incertain.
\newcommand{\uncertain}[1]{\uline{#1}}

% Ajout de l'éditeur — une lettre manquante, un mot sous-entendu.
\newcommand{\add}[1]{\textcolor{blue!50!black}{[#1]}}

% Biffé par Grothendieck lui-même. Ce qu'il a rayé fait partie du document :
% une rature dit souvent où la pensée a changé de direction.
\newcommand{\struck}[1]{\sout{#1}}

% Note du transcripteur, sur l'état matériel du feuillet.
\newcommand{\note}[1]{\par\smallskip{\small\color{gray!60!black}\textit{[#1]}}\par\smallskip}

% Note marginale de Grothendieck, distincte du corps du texte.
\newcommand{\marginal}[1]{\par\smallskip\noindent\hspace{1em}%
  {\small\color{gray!60!black}#1}\par\smallskip}

% --- Titre ------------------------------------------------------------------

\AtBeginDocument{%
  \begin{center}
    {\large\bfseries Fonds Alexandre Grothendieck --- cote n\textsuperscript{o}~\grothFolder}\\[2pt]
    {\itshape\grothFoldertitle}\\[4pt]
    {\small feuillets \grothFirst--\grothLast\ (lot \grothBatch)}\\[2pt]
    {\footnotesize\color{gray!60!black}Transcription établie d'après le fac-similé numérisé
     par l'Université de Montpellier.\\
     Les lectures incertaines sont soulignées, les illisibles notées [\,\ldots\,],
     les ajouts entre crochets.}
  \end{center}
  \vspace{1em}}

\endinput


\folder{116}
\batch{1}
\pages{1}{10}
\dating{[à partir de 1981]}
\watermark{Édition de démonstration}
\foldertitle{Mapping-cône etc. : notes manuscrites (s.d.)}

\begin{document}

\section*{Mapping cone}

\note{Ces deux mots sont de sa main, seuls en haut à droite du premier
feuillet, une chemise de papier kraft. Le feuillet ne porte rien d'autre et ne
prend pas de \texttt{page} ; c'est de là que vient le titre du dossier, qui
n'est donc pas des archivistes. La page 2 est l'autre face de la même chemise.}

\page{2}

\note{Des notations jetées, sans phrase suivie : le triple adjoint d'une
immersion, les deux catégories de préfaisceaux $\widehat{A}_0$ et
$\widehat{A}$, et un petit schéma de flèches. La page se lit comme un
pense-bête, non comme un raisonnement.}

\[
\uncertain{{}_{b}\backslash A} \longleftarrow A_{b}
\]

\begin{gather*}
i_{!} = \ill{} \\
i^{*} = \ill{} \\
i_{*}
\end{gather*}

\note{Les trois lignes sont écrites l'une sous l'autre dans le coin supérieur
droit. Le membre de droite des deux premières est un même glyphe, coupé en
haut par le bord du feuillet sur la première ; il n'est pas lu ici.}

\[
p_{A} : A \longrightarrow A_{0}
\]

\begin{tikzcd}[row sep=large, column sep=large]
b_0 \arrow[r, "p"] \arrow[d] & b' \\
b_0 \arrow[ur, "p"'] &
\end{tikzcd}

\note{Sur le feuillet, les deux $b_0$ sont en outre joints par un double trait
vertical. Du haut de la colonne de gauche part une flèche en gras vers un
objet non nommé, marqué d'une croix, d'où redescend une flèche verticale ;
faute de pouvoir nommer ce but, ces deux flèches ne sont pas reproduites.}

\[
\begin{array}{c}
p_{A!}(F) \\
\uparrow \\
p_{A!}\, i_{!}\bigl(\ill{}(F)\bigr)
\end{array}
\]

\note{La flèche verticale est de lui et va de la ligne du bas vers celle du
haut ; l'argument de $i_{!}$ n'est pas lu.}

\begin{gather*}
i_{!}\, i^{*}(F) \xrightarrow{\ \sim\ } F \\
i_{*}\, i^{*} F \xleftarrow{\ \sim\ } F
\end{gather*}

\begin{tikzcd}[column sep=huge]
\widehat{A}_0 \arrow[r, "i_{!}", bend left=25] \arrow[r, "i_{*}"', bend right=25] & \widehat{A} \arrow[l, "i^{0}" description]
\end{tikzcd}

\note{Les trois flèches sont superposées entre les deux mêmes objets : $i_{!}$
en haut et $i_{*}$ en bas vers la droite, celle du milieu vers la gauche.}

\begin{itemize}
  \item $i_0^{*}$ \uncertain{pl. fid.}
  \item $i_{0*}$ localisation
  \item $i_{0!}$ localisation
\end{itemize}

\note{Les trois lignes sont prises dans une accolade de sa main. Les indices
$0$ sont écrits en exposant sur le $i$, et l'abréviation de la première ligne
n'est pas assurée.}

\[
F \longrightarrow \ill{}\ \ill{}\, F \quad \text{induit iso sur } \Gamma, \text{ puis } H^{*} F
\]

\[
\ill{} \quad \mathrm{Ex}_{\alpha} \quad \ill{} \qquad\qquad A_0 \longrightarrow A
\]

\section*{Relations entre morphismes de topos et recollements de topos}

\note{Titre de sa main, en tête de la page 3. Dans tout ce qui suit
$\mathcal{X}$ désigne un topos et $\mathcal{F}(\mathcal{X})$ la catégorie de
ses faisceaux ; les lettres droites $X$, $X_0$, $X_1$ des pages 6 et 7
désignent des catégories, et $\widehat{X}$ la catégorie des préfaisceaux sur
$X$. La distinction typographique est de l'édition ; le feuillet ne sépare les
deux tracés que par la forme de la lettre.}

\page{3}

1) \struck{Soit} Soit
\[
(1) \qquad f : \mathcal{X}_0 \longrightarrow \mathcal{X}_1
\]
un morphisme de topos, on peut le considérer comme topos fibré sur $\Delta_1$
(fibres $\mathcal{X}_0$, $\mathcal{X}_1$), et regarder le topos total
\uncertain{associé} $\mathcal{X} = C(f)$ (le mapping-cone) dont les faisceaux
sont les systèmes
\[
(2) \qquad
\begin{array}{l}
(F_0, F_1, u), \qquad F_0 \in \mathrm{Ob}\,\mathcal{F}(\mathcal{X}_0),\ F_1 \in \mathrm{Ob}\,\mathcal{F}(\mathcal{X}_1) \\[4pt]
\qquad u : f^{*}(F_1) \longrightarrow F_0 \quad \text{i.e.} \\[4pt]
\qquad F_1 \longrightarrow f_{*}(F_0)
\end{array}
\]

Alors $\mathcal{X}_0$ peut être regardé comme l'ouvert de $\mathcal{X}$,
$\mathcal{X}_1$ comme le fermé complémentaire, mais de plus l'inclusion
\[
(3) \qquad \mathcal{X}_1 \xrightarrow{\ i_1\ } \mathcal{X}
\]
est strictement \struck{asphérique} \add{équiv. f.}, i.e. pour tout faisceau
$F$ sur $\mathcal{X}$, \add{$F = (F_0, F_1, u)$}
\[
H^{*}(\mathcal{X}, F) \xrightarrow{\ \sim\ } H^{*}\bigl(\mathcal{X}_1, i_1^{*}(F)\bigr)
\]

\note{Les mots « équiv. f. » sont écrits au-dessus de « asphérique » barré, et
« est strictement » est souligné d'un trait ondulé. La même substitution
revient en tête de la page 5.}

Le « foncteur véritablement \ill{} » \uncertain{(avec i.g.)} (de la théorie
\ill{})
\[
\varphi : \mathcal{F}(\mathcal{X}_0) \longrightarrow \mathcal{F}(\mathcal{X}_1)
\]
n'est autre que
\[
(4) \qquad \varphi = f_{*}
\]

\marginal{N.B. On a \ill{} $\mathcal{X} \longrightarrow \mathrm{Top}(\Delta_1)$,
on a \ill{} qui revient au $=$, Alors \ill{} on a des $\mathcal{X}_0$ \ill{}
(i.e. \ill{} $0$ de l'objet \ill{}) topos ouvert \ill{} $\{0\}$ de \ill{} du
sous-objet \ill{} final de $\widehat{\Delta_1}$) 3)}

\marginal{Re N.B. On a une rétraction \ill{} $\mathcal{X}$ sur le $2^{\text{ème}}$
\ill{} $\mathcal{X}_1$, par le foncteur \ill{} (topos.)
$i_{*}(F_1) = (f^{*}(F_1), F_1, \uncertain{\mathrm{id}})$,
$\mathcal{X} = (\mathcal{X}_0 \times \widehat{\Delta_1}) \amalg_{(\mathcal{X}_0, 1)} \mathcal{X}_1$}

\note{Ces deux blocs sont écrits en biais, à quelque 40 ou 45°, le long du
bord gauche du feuillet, dans une main plus rapide que celle du corps ; ils
sont donnés aussi loin qu'ils vont.}

\page{4}

Donc : se donner un morphisme de topos \add{$\mathcal{X}_0 \to \mathcal{X}_1$}
revient au $=$ que de se donner un topos $\mathcal{X}$, (un sous-topos ouvert
$\mathcal{X}_0$) et le fermé complémentaire $\mathcal{X}_1$, de telle façon,
que le foncteur \ill{} $\varphi$ soit, non seulement exact à g\add{auche},
mais commutant aux $\varprojlim$ quelconques, \struck{\ill{}} i.e. admettant
un adjoint à gauche $f^{*}$, et que de plus ce dernier soit exact à gauche.

\note{Il écrit le signe $=$ pour le mot « même » ; il le refait à la page 6.
Le mot qui qualifie ici le foncteur $\varphi$ n'est pas lu : le tracé
ressemble à celui de « véritablement » à la page précédente, mais la syntaxe
ne le prend pas.}

Alors l'inclusion $i : \mathcal{X}_1 \longrightarrow \mathcal{X}$ est strict.
asphérique (réciproque ?) et la suite exacte de cohomologie relative

\marginal{?}

\[
(5) \qquad \longrightarrow H^{i}_{\mathcal{X}_1}(\mathcal{X}, F) \longrightarrow H^{i}(\mathcal{X}, F) \longrightarrow H^{i}(\mathcal{X}_0, F) \longrightarrow H^{i+1}_{\mathcal{X}_1}(\mathcal{X}, F) \longrightarrow \cdots
\]

peut s'écrire

\[
(5') \qquad \longrightarrow H^{i}_{\mathcal{X}_1}(\mathcal{X}, F) \longrightarrow H^{i}(\mathcal{X}_1, F_1) \longrightarrow H^{i}(\mathcal{X}_0, F_0) \longrightarrow H^{i+1}_{\mathcal{X}_1}(\mathcal{X}, F) \longrightarrow \cdots
\]

On pose
\[
H^{i}_{\mathcal{X}_1}(\mathcal{X}, F) \overset{\text{déf}}{=} H^{i}\bigl(\mathcal{X}_1 \bmod \mathcal{X}_0,\ F = (F_0, F_1, u)\bigr)
\]
\note{Sous « $\mathcal{X}_1 \bmod \mathcal{X}_0$ » il porte une accolade et,
dessous, « $f$ sous-entendu ».}
\ill{} aussi \struck{\ill{}}
\[
H^{i}\bigl(\{\mathcal{X}_1 \bmod \mathcal{X}_0\},\ F_1\bigr) \qquad \text{si } F_0 = f^{*}(F_1),\ u \text{ tautolog.}
\]

\page{5}

Dire que $f : \mathcal{X}_0 \longrightarrow \mathcal{X}_1$ est
\struck{asphérique} \add{équiv. faible}, signifie aussi que
$\mathcal{X}_0 \longrightarrow \mathcal{X}$ (immersion ouverte) l'est.

N.B. Les faisceaux loc. constants sur $\mathcal{X}$ sont ceux de la forme
$(F_1, f^{*}(F_1), u\ \text{tautol.})$, avec $F_1$ loc. \uncertain{cst.} sur
$\mathcal{X}_1$.

\note{L'ordre des deux premiers termes est inversé par rapport à la
convention $(F_0, F_1, u)$ posée à la page 3 ; le feuillet est laissé tel
quel. Un faisceau de flèches part du mot $\mathcal{X}$ vers la ligne
précédente, où il corrige un premier $\mathcal{X}_1$.}

2) Partons inversement d'un topos $\mathcal{X}$ et d'un sous-topos ouvert
\[
(6) \qquad \mathcal{X}_0 \subset \mathcal{X},
\]
$\Longleftrightarrow$ sous-objet de l'objet final $e_{\mathcal{X}}$ de
$\mathcal{F}(\mathcal{X})$, d'où topos résiduel $\mathcal{X}_1$. On peut
considérer aussi que la donnée de $\mathcal{X}_0$ correspond \ill{} celle
d'un morphisme
\[
\pi : \mathcal{X} \longrightarrow \mathrm{Top}(\Delta_1)
\]
$\mathcal{X}_0$ étant l'image inverse du sous-topos ouvert « universel »
$\{0\}$ dans $\mathrm{Top}(\Delta_1)$, et $\mathcal{X}_1$ \struck{\ill{}}
l'image inverse du sous-topos fermé $\{1\}$.

\page{6}

D'autre part, la donnée de $(\mathcal{X}, \mathcal{X}_0)$ équivaut à celle de
\[
(\mathcal{X}_0, \mathcal{X}_1, \varphi) \qquad
\varphi : \mathcal{F}(\mathcal{X}_0) \longrightarrow \mathcal{F}(\mathcal{X}_1)
\]
\struck{foncteur} exact : g.

\note{Au-dessus du $\varphi$ il inscrit, entre crochets, $f^{*}$ surmontant un
double trait vertical.}

On s'intéresse à des conditions \add{« se recolle » sur $(\mathcal{X}_0, \mathcal{X}_1, \varphi)$}
qui assurent que les \add{deux} inclusions
\[
\mathcal{X}_0 \longrightarrow \mathcal{X} \longleftarrow \mathcal{X}_1
\]
soient des équivalences faibles. Les cas qui nous intéressent le plus, sont
les cas où $\varphi = f_{*}$ commutant \add{en fait} aux $\varprojlim$
quelconques, i.e. \ill{} un adjoint à droite $f^{*}$ — comme p. ex. celui du
§ précédent, où on suppose de plus que $f^{*}$ est exact à gauche (mais on n'a
pas envie de le supposer ici).

\note{« à droite » : les pages 4 et 7 écrivent « à gauche » au même endroit de
la même phrase. Le feuillet est laissé tel quel.}

Un cas particulièrement important pour nous est celui où $\mathcal{X}$ est de
la forme \struck{\ill{}} \add{$\mathrm{Top}(X)$}, \ill{} $X$ dans
$(\mathrm{Cat})$, il revient au $=$ de dire que

\page{7}

$\mathcal{X}_0$, $\mathcal{X}_1$ sont de la forme
\struck{$\widehat{X}_0$, $\widehat{X}_1$}
\add{$\mathrm{Top}(X_0)$, $\mathrm{Top}(X_1)$}, et que $\varphi$ commute aux
$\varprojlim$ quelconques i.e. \ill{} un adjoint à g. $f^{*}$, qui est bien
défini par un foncteur
\[
X_1 \longrightarrow \widehat{X}_0
\]
(\ill{} par « continuité » en $\varinjlim$
\[
\widehat{X}_1 \xrightarrow{\ f^{*}\ } \widehat{X}_0
\]
), correspondant i.e. à un bifoncteur
\[
(x_0, x_1) \mapsto h(x_1, x_0), \qquad X_0^{\mathrm{op}} \times X_1 \longrightarrow (\mathrm{Ens}).
\]

\section*{Feuillet séparé}

\note{La page 9 n'est pas sur le sujet des pages 3 à 7 : elle porte, sur le
tiers supérieur d'un feuillet par ailleurs vide, un carré de flèches et une
comparaison de deux limites projectives sur une catégorie filtrante.}

\page{9}

\begin{tikzcd}
A \arrow[d] & A' \arrow[l] \arrow[d] \\
B & B' \arrow[l]
\end{tikzcd}

$A$ filtrant
\[
A_{/b} \longleftrightarrow A_{b}
\]
\[
\varprojlim_{A_{/b}} \bigl(F \mid A_{/b}\bigr) \overset{\text{iso}}{\Longrightarrow} \varprojlim_{A_b} F \mid A_{b}
\]
Pour tout $F$ sur $A$ (ou sur $A_{/b}$)

\marginal{$f$ dans \uncertain{Hot}}

\note{Ces trois derniers mots sont écrits verticalement le long du bord droit
du feuillet.}

\end{document}
